\chapter{Floating Point and the Epsilon Bug}\label{ch:float}

\begin{objectives}
\begin{itemize}
\item read an IEEE-754 binary64 number field by field, and give the exact stored values of $0.9$ and $0.1$;
\item explain, with Sterbenz's lemma, why \cc{1.0 - 0.90} is \emph{exactly} $0.09999999999999997779\ldots$;
\item trace how truncation turned \cc{tail_index} 1 into 0 at $n=10$, $c=0.90$;
\item prove a sufficient condition under which the repository's $\varepsilon=10^{-9}$ guard is correct, and exhibit an input where it is not;
\item compare the alternatives: integer basis points, rounding with tolerance, exact rationals.
\end{itemize}
{\small\color{muted}\textbf{Chapter references:} IEEE 754-2019~\cite{ieee754-2019}; Goldberg~\cite{goldberg1991what}; Higham~\cite{higham2002accuracy};
Sterbenz~\cite{sterbenz1974floating}.}
\end{objectives}

\section{The bug, as it was found}
While spot-checking \cc{historical_var} against hand-computed quantiles on the ten-point series of
\cref{ex:tenpoint}, the exact-boundary case $n=10$, $c=0.90$ produced \cc{tail_index} $=0$, although one tenth of ten
is one, so the function reported ``not enough returns''. The truncating cast had no epsilon (\file{CHANGELOG.md},
entry of 2026-10-17). The repository now reads:

\repo[firstline=32,lastline=42]{src/statistics/risk.cpp}{The fix and its comment, verbatim. The rest of this chapter asks
why the comment is true, and when the fix is enough.}{lst:epsfix}

\begin{progout}{risk\_check: the tail fraction and its truncation, with and without epsilon}{lst:epsout}
c=0.95 1-c=0.050000000000000044 (1-c)*10=0.50000000000000044 floor_noeps=0 floor_eps=0
c=0.90 1-c=0.099999999999999978 (1-c)*10=0.99999999999999978 floor_noeps=0 floor_eps=1
c=0.80 1-c=0.19999999999999996 (1-c)*10=1.9999999999999996 floor_noeps=1 floor_eps=2
\end{progout}

The bug was not confined to $0.90$: at $c=0.80$ the unguarded code would have used one tail observation instead of two and
reported $\VaR=0.05$ instead of $0.04$, a silently wrong answer rather than an exception.

\section{IEEE-754 binary64}
\begin{definition}[binary64~\cite{ieee754-2019}]\label[definition]{def:binary64}
A 64-bit word with sign bit $s$, 11-bit biased exponent $E$ and 52-bit fraction $f$ represents, for $0<E<2047$,
\[ x=(-1)^s\cdot\Bigl(1+\frac{f}{2^{52}}\Bigr)\cdot 2^{E-1023}. \]
Every such $x$ is a dyadic rational: an integer times a power of two.
\end{definition}

\begin{lemma}[Which decimals are exact]\label[lemma]{lem:dyadic}
A rational $a/b$ in lowest terms is a binary64 value (within range) only if $b$ is a power of two. In particular
$0.9=9/10$, $0.1=1/10$, $0.05=1/20$ and $0.95=19/20$ are not representable; $0.5$, $0.25$ and $0.75$ are.
\end{lemma}
\begin{proof}
By \cref{def:binary64}, $x=M\cdot2^{e}$ with integers $M$, $e$. If $e\ge0$, $x$ is an integer; otherwise
$x=M/2^{-e}$, whose reduced denominator divides $2^{-e}$. The denominators $10$ and $20$ contain the factor $5$.
\end{proof}

\begin{figure}[htbp]
\centering
\resizebox{\linewidth}{!}{% generated by gen_bits.py from struct.pack of the binary64 values
\begin{tikzpicture}[x=0.215cm,y=0.62cm]
\node[anchor=east,font=\small\ttfamily,text=fg] at (-0.6,0.25) {0.9};
\node[anchor=west,font=\scriptsize\ttfamily,text=org] at (64.6,0.25) {0x3FECCCCCCCCCCCCD};
\draw[draw=red!70!bg,fill=codebg] (0,0.0) rectangle ++(1,0.5); \node[font=\tiny\ttfamily,text=fg] at (0.5,0.25) {0};
\draw[draw=yel!70!bg,fill=codebg] (1,0.0) rectangle ++(1,0.5); \node[font=\tiny\ttfamily,text=fg] at (1.5,0.25) {0};
\draw[draw=yel!70!bg,fill=yel!45!bg] (2,0.0) rectangle ++(1,0.5); \node[font=\tiny\ttfamily,text=fg] at (2.5,0.25) {1};
\draw[draw=yel!70!bg,fill=yel!45!bg] (3,0.0) rectangle ++(1,0.5); \node[font=\tiny\ttfamily,text=fg] at (3.5,0.25) {1};
\draw[draw=yel!70!bg,fill=yel!45!bg] (4,0.0) rectangle ++(1,0.5); \node[font=\tiny\ttfamily,text=fg] at (4.5,0.25) {1};
\draw[draw=yel!70!bg,fill=yel!45!bg] (5,0.0) rectangle ++(1,0.5); \node[font=\tiny\ttfamily,text=fg] at (5.5,0.25) {1};
\draw[draw=yel!70!bg,fill=yel!45!bg] (6,0.0) rectangle ++(1,0.5); \node[font=\tiny\ttfamily,text=fg] at (6.5,0.25) {1};
\draw[draw=yel!70!bg,fill=yel!45!bg] (7,0.0) rectangle ++(1,0.5); \node[font=\tiny\ttfamily,text=fg] at (7.5,0.25) {1};
\draw[draw=yel!70!bg,fill=yel!45!bg] (8,0.0) rectangle ++(1,0.5); \node[font=\tiny\ttfamily,text=fg] at (8.5,0.25) {1};
\draw[draw=yel!70!bg,fill=yel!45!bg] (9,0.0) rectangle ++(1,0.5); \node[font=\tiny\ttfamily,text=fg] at (9.5,0.25) {1};
\draw[draw=yel!70!bg,fill=yel!45!bg] (10,0.0) rectangle ++(1,0.5); \node[font=\tiny\ttfamily,text=fg] at (10.5,0.25) {1};
\draw[draw=yel!70!bg,fill=codebg] (11,0.0) rectangle ++(1,0.5); \node[font=\tiny\ttfamily,text=fg] at (11.5,0.25) {0};
\draw[draw=acc!70!bg,fill=acc!45!bg] (12,0.0) rectangle ++(1,0.5); \node[font=\tiny\ttfamily,text=fg] at (12.5,0.25) {1};
\draw[draw=acc!70!bg,fill=acc!45!bg] (13,0.0) rectangle ++(1,0.5); \node[font=\tiny\ttfamily,text=fg] at (13.5,0.25) {1};
\draw[draw=acc!70!bg,fill=codebg] (14,0.0) rectangle ++(1,0.5); \node[font=\tiny\ttfamily,text=fg] at (14.5,0.25) {0};
\draw[draw=acc!70!bg,fill=codebg] (15,0.0) rectangle ++(1,0.5); \node[font=\tiny\ttfamily,text=fg] at (15.5,0.25) {0};
\draw[draw=acc!70!bg,fill=acc!45!bg] (16,0.0) rectangle ++(1,0.5); \node[font=\tiny\ttfamily,text=fg] at (16.5,0.25) {1};
\draw[draw=acc!70!bg,fill=acc!45!bg] (17,0.0) rectangle ++(1,0.5); \node[font=\tiny\ttfamily,text=fg] at (17.5,0.25) {1};
\draw[draw=acc!70!bg,fill=codebg] (18,0.0) rectangle ++(1,0.5); \node[font=\tiny\ttfamily,text=fg] at (18.5,0.25) {0};
\draw[draw=acc!70!bg,fill=codebg] (19,0.0) rectangle ++(1,0.5); \node[font=\tiny\ttfamily,text=fg] at (19.5,0.25) {0};
\draw[draw=acc!70!bg,fill=acc!45!bg] (20,0.0) rectangle ++(1,0.5); \node[font=\tiny\ttfamily,text=fg] at (20.5,0.25) {1};
\draw[draw=acc!70!bg,fill=acc!45!bg] (21,0.0) rectangle ++(1,0.5); \node[font=\tiny\ttfamily,text=fg] at (21.5,0.25) {1};
\draw[draw=acc!70!bg,fill=codebg] (22,0.0) rectangle ++(1,0.5); \node[font=\tiny\ttfamily,text=fg] at (22.5,0.25) {0};
\draw[draw=acc!70!bg,fill=codebg] (23,0.0) rectangle ++(1,0.5); \node[font=\tiny\ttfamily,text=fg] at (23.5,0.25) {0};
\draw[draw=acc!70!bg,fill=acc!45!bg] (24,0.0) rectangle ++(1,0.5); \node[font=\tiny\ttfamily,text=fg] at (24.5,0.25) {1};
\draw[draw=acc!70!bg,fill=acc!45!bg] (25,0.0) rectangle ++(1,0.5); \node[font=\tiny\ttfamily,text=fg] at (25.5,0.25) {1};
\draw[draw=acc!70!bg,fill=codebg] (26,0.0) rectangle ++(1,0.5); \node[font=\tiny\ttfamily,text=fg] at (26.5,0.25) {0};
\draw[draw=acc!70!bg,fill=codebg] (27,0.0) rectangle ++(1,0.5); \node[font=\tiny\ttfamily,text=fg] at (27.5,0.25) {0};
\draw[draw=acc!70!bg,fill=acc!45!bg] (28,0.0) rectangle ++(1,0.5); \node[font=\tiny\ttfamily,text=fg] at (28.5,0.25) {1};
\draw[draw=acc!70!bg,fill=acc!45!bg] (29,0.0) rectangle ++(1,0.5); \node[font=\tiny\ttfamily,text=fg] at (29.5,0.25) {1};
\draw[draw=acc!70!bg,fill=codebg] (30,0.0) rectangle ++(1,0.5); \node[font=\tiny\ttfamily,text=fg] at (30.5,0.25) {0};
\draw[draw=acc!70!bg,fill=codebg] (31,0.0) rectangle ++(1,0.5); \node[font=\tiny\ttfamily,text=fg] at (31.5,0.25) {0};
\draw[draw=acc!70!bg,fill=acc!45!bg] (32,0.0) rectangle ++(1,0.5); \node[font=\tiny\ttfamily,text=fg] at (32.5,0.25) {1};
\draw[draw=acc!70!bg,fill=acc!45!bg] (33,0.0) rectangle ++(1,0.5); \node[font=\tiny\ttfamily,text=fg] at (33.5,0.25) {1};
\draw[draw=acc!70!bg,fill=codebg] (34,0.0) rectangle ++(1,0.5); \node[font=\tiny\ttfamily,text=fg] at (34.5,0.25) {0};
\draw[draw=acc!70!bg,fill=codebg] (35,0.0) rectangle ++(1,0.5); \node[font=\tiny\ttfamily,text=fg] at (35.5,0.25) {0};
\draw[draw=acc!70!bg,fill=acc!45!bg] (36,0.0) rectangle ++(1,0.5); \node[font=\tiny\ttfamily,text=fg] at (36.5,0.25) {1};
\draw[draw=acc!70!bg,fill=acc!45!bg] (37,0.0) rectangle ++(1,0.5); \node[font=\tiny\ttfamily,text=fg] at (37.5,0.25) {1};
\draw[draw=acc!70!bg,fill=codebg] (38,0.0) rectangle ++(1,0.5); \node[font=\tiny\ttfamily,text=fg] at (38.5,0.25) {0};
\draw[draw=acc!70!bg,fill=codebg] (39,0.0) rectangle ++(1,0.5); \node[font=\tiny\ttfamily,text=fg] at (39.5,0.25) {0};
\draw[draw=acc!70!bg,fill=acc!45!bg] (40,0.0) rectangle ++(1,0.5); \node[font=\tiny\ttfamily,text=fg] at (40.5,0.25) {1};
\draw[draw=acc!70!bg,fill=acc!45!bg] (41,0.0) rectangle ++(1,0.5); \node[font=\tiny\ttfamily,text=fg] at (41.5,0.25) {1};
\draw[draw=acc!70!bg,fill=codebg] (42,0.0) rectangle ++(1,0.5); \node[font=\tiny\ttfamily,text=fg] at (42.5,0.25) {0};
\draw[draw=acc!70!bg,fill=codebg] (43,0.0) rectangle ++(1,0.5); \node[font=\tiny\ttfamily,text=fg] at (43.5,0.25) {0};
\draw[draw=acc!70!bg,fill=acc!45!bg] (44,0.0) rectangle ++(1,0.5); \node[font=\tiny\ttfamily,text=fg] at (44.5,0.25) {1};
\draw[draw=acc!70!bg,fill=acc!45!bg] (45,0.0) rectangle ++(1,0.5); \node[font=\tiny\ttfamily,text=fg] at (45.5,0.25) {1};
\draw[draw=acc!70!bg,fill=codebg] (46,0.0) rectangle ++(1,0.5); \node[font=\tiny\ttfamily,text=fg] at (46.5,0.25) {0};
\draw[draw=acc!70!bg,fill=codebg] (47,0.0) rectangle ++(1,0.5); \node[font=\tiny\ttfamily,text=fg] at (47.5,0.25) {0};
\draw[draw=acc!70!bg,fill=acc!45!bg] (48,0.0) rectangle ++(1,0.5); \node[font=\tiny\ttfamily,text=fg] at (48.5,0.25) {1};
\draw[draw=acc!70!bg,fill=acc!45!bg] (49,0.0) rectangle ++(1,0.5); \node[font=\tiny\ttfamily,text=fg] at (49.5,0.25) {1};
\draw[draw=acc!70!bg,fill=codebg] (50,0.0) rectangle ++(1,0.5); \node[font=\tiny\ttfamily,text=fg] at (50.5,0.25) {0};
\draw[draw=acc!70!bg,fill=codebg] (51,0.0) rectangle ++(1,0.5); \node[font=\tiny\ttfamily,text=fg] at (51.5,0.25) {0};
\draw[draw=acc!70!bg,fill=acc!45!bg] (52,0.0) rectangle ++(1,0.5); \node[font=\tiny\ttfamily,text=fg] at (52.5,0.25) {1};
\draw[draw=acc!70!bg,fill=acc!45!bg] (53,0.0) rectangle ++(1,0.5); \node[font=\tiny\ttfamily,text=fg] at (53.5,0.25) {1};
\draw[draw=acc!70!bg,fill=codebg] (54,0.0) rectangle ++(1,0.5); \node[font=\tiny\ttfamily,text=fg] at (54.5,0.25) {0};
\draw[draw=acc!70!bg,fill=codebg] (55,0.0) rectangle ++(1,0.5); \node[font=\tiny\ttfamily,text=fg] at (55.5,0.25) {0};
\draw[draw=acc!70!bg,fill=acc!45!bg] (56,0.0) rectangle ++(1,0.5); \node[font=\tiny\ttfamily,text=fg] at (56.5,0.25) {1};
\draw[draw=acc!70!bg,fill=acc!45!bg] (57,0.0) rectangle ++(1,0.5); \node[font=\tiny\ttfamily,text=fg] at (57.5,0.25) {1};
\draw[draw=acc!70!bg,fill=codebg] (58,0.0) rectangle ++(1,0.5); \node[font=\tiny\ttfamily,text=fg] at (58.5,0.25) {0};
\draw[draw=acc!70!bg,fill=codebg] (59,0.0) rectangle ++(1,0.5); \node[font=\tiny\ttfamily,text=fg] at (59.5,0.25) {0};
\draw[draw=acc!70!bg,fill=acc!45!bg] (60,0.0) rectangle ++(1,0.5); \node[font=\tiny\ttfamily,text=fg] at (60.5,0.25) {1};
\draw[draw=acc!70!bg,fill=acc!45!bg] (61,0.0) rectangle ++(1,0.5); \node[font=\tiny\ttfamily,text=fg] at (61.5,0.25) {1};
\draw[draw=acc!70!bg,fill=codebg] (62,0.0) rectangle ++(1,0.5); \node[font=\tiny\ttfamily,text=fg] at (62.5,0.25) {0};
\draw[draw=acc!70!bg,fill=acc!45!bg] (63,0.0) rectangle ++(1,0.5); \node[font=\tiny\ttfamily,text=fg] at (63.5,0.25) {1};
\node[anchor=east,font=\small\ttfamily,text=fg] at (-0.6,-1.30) {0.1};
\node[anchor=west,font=\scriptsize\ttfamily,text=org] at (64.6,-1.30) {0x3FB999999999999A};
\draw[draw=red!70!bg,fill=codebg] (0,-1.55) rectangle ++(1,0.5); \node[font=\tiny\ttfamily,text=fg] at (0.5,-1.30) {0};
\draw[draw=yel!70!bg,fill=codebg] (1,-1.55) rectangle ++(1,0.5); \node[font=\tiny\ttfamily,text=fg] at (1.5,-1.30) {0};
\draw[draw=yel!70!bg,fill=yel!45!bg] (2,-1.55) rectangle ++(1,0.5); \node[font=\tiny\ttfamily,text=fg] at (2.5,-1.30) {1};
\draw[draw=yel!70!bg,fill=yel!45!bg] (3,-1.55) rectangle ++(1,0.5); \node[font=\tiny\ttfamily,text=fg] at (3.5,-1.30) {1};
\draw[draw=yel!70!bg,fill=yel!45!bg] (4,-1.55) rectangle ++(1,0.5); \node[font=\tiny\ttfamily,text=fg] at (4.5,-1.30) {1};
\draw[draw=yel!70!bg,fill=yel!45!bg] (5,-1.55) rectangle ++(1,0.5); \node[font=\tiny\ttfamily,text=fg] at (5.5,-1.30) {1};
\draw[draw=yel!70!bg,fill=yel!45!bg] (6,-1.55) rectangle ++(1,0.5); \node[font=\tiny\ttfamily,text=fg] at (6.5,-1.30) {1};
\draw[draw=yel!70!bg,fill=yel!45!bg] (7,-1.55) rectangle ++(1,0.5); \node[font=\tiny\ttfamily,text=fg] at (7.5,-1.30) {1};
\draw[draw=yel!70!bg,fill=yel!45!bg] (8,-1.55) rectangle ++(1,0.5); \node[font=\tiny\ttfamily,text=fg] at (8.5,-1.30) {1};
\draw[draw=yel!70!bg,fill=codebg] (9,-1.55) rectangle ++(1,0.5); \node[font=\tiny\ttfamily,text=fg] at (9.5,-1.30) {0};
\draw[draw=yel!70!bg,fill=yel!45!bg] (10,-1.55) rectangle ++(1,0.5); \node[font=\tiny\ttfamily,text=fg] at (10.5,-1.30) {1};
\draw[draw=yel!70!bg,fill=yel!45!bg] (11,-1.55) rectangle ++(1,0.5); \node[font=\tiny\ttfamily,text=fg] at (11.5,-1.30) {1};
\draw[draw=acc!70!bg,fill=acc!45!bg] (12,-1.55) rectangle ++(1,0.5); \node[font=\tiny\ttfamily,text=fg] at (12.5,-1.30) {1};
\draw[draw=acc!70!bg,fill=codebg] (13,-1.55) rectangle ++(1,0.5); \node[font=\tiny\ttfamily,text=fg] at (13.5,-1.30) {0};
\draw[draw=acc!70!bg,fill=codebg] (14,-1.55) rectangle ++(1,0.5); \node[font=\tiny\ttfamily,text=fg] at (14.5,-1.30) {0};
\draw[draw=acc!70!bg,fill=acc!45!bg] (15,-1.55) rectangle ++(1,0.5); \node[font=\tiny\ttfamily,text=fg] at (15.5,-1.30) {1};
\draw[draw=acc!70!bg,fill=acc!45!bg] (16,-1.55) rectangle ++(1,0.5); \node[font=\tiny\ttfamily,text=fg] at (16.5,-1.30) {1};
\draw[draw=acc!70!bg,fill=codebg] (17,-1.55) rectangle ++(1,0.5); \node[font=\tiny\ttfamily,text=fg] at (17.5,-1.30) {0};
\draw[draw=acc!70!bg,fill=codebg] (18,-1.55) rectangle ++(1,0.5); \node[font=\tiny\ttfamily,text=fg] at (18.5,-1.30) {0};
\draw[draw=acc!70!bg,fill=acc!45!bg] (19,-1.55) rectangle ++(1,0.5); \node[font=\tiny\ttfamily,text=fg] at (19.5,-1.30) {1};
\draw[draw=acc!70!bg,fill=acc!45!bg] (20,-1.55) rectangle ++(1,0.5); \node[font=\tiny\ttfamily,text=fg] at (20.5,-1.30) {1};
\draw[draw=acc!70!bg,fill=codebg] (21,-1.55) rectangle ++(1,0.5); \node[font=\tiny\ttfamily,text=fg] at (21.5,-1.30) {0};
\draw[draw=acc!70!bg,fill=codebg] (22,-1.55) rectangle ++(1,0.5); \node[font=\tiny\ttfamily,text=fg] at (22.5,-1.30) {0};
\draw[draw=acc!70!bg,fill=acc!45!bg] (23,-1.55) rectangle ++(1,0.5); \node[font=\tiny\ttfamily,text=fg] at (23.5,-1.30) {1};
\draw[draw=acc!70!bg,fill=acc!45!bg] (24,-1.55) rectangle ++(1,0.5); \node[font=\tiny\ttfamily,text=fg] at (24.5,-1.30) {1};
\draw[draw=acc!70!bg,fill=codebg] (25,-1.55) rectangle ++(1,0.5); \node[font=\tiny\ttfamily,text=fg] at (25.5,-1.30) {0};
\draw[draw=acc!70!bg,fill=codebg] (26,-1.55) rectangle ++(1,0.5); \node[font=\tiny\ttfamily,text=fg] at (26.5,-1.30) {0};
\draw[draw=acc!70!bg,fill=acc!45!bg] (27,-1.55) rectangle ++(1,0.5); \node[font=\tiny\ttfamily,text=fg] at (27.5,-1.30) {1};
\draw[draw=acc!70!bg,fill=acc!45!bg] (28,-1.55) rectangle ++(1,0.5); \node[font=\tiny\ttfamily,text=fg] at (28.5,-1.30) {1};
\draw[draw=acc!70!bg,fill=codebg] (29,-1.55) rectangle ++(1,0.5); \node[font=\tiny\ttfamily,text=fg] at (29.5,-1.30) {0};
\draw[draw=acc!70!bg,fill=codebg] (30,-1.55) rectangle ++(1,0.5); \node[font=\tiny\ttfamily,text=fg] at (30.5,-1.30) {0};
\draw[draw=acc!70!bg,fill=acc!45!bg] (31,-1.55) rectangle ++(1,0.5); \node[font=\tiny\ttfamily,text=fg] at (31.5,-1.30) {1};
\draw[draw=acc!70!bg,fill=acc!45!bg] (32,-1.55) rectangle ++(1,0.5); \node[font=\tiny\ttfamily,text=fg] at (32.5,-1.30) {1};
\draw[draw=acc!70!bg,fill=codebg] (33,-1.55) rectangle ++(1,0.5); \node[font=\tiny\ttfamily,text=fg] at (33.5,-1.30) {0};
\draw[draw=acc!70!bg,fill=codebg] (34,-1.55) rectangle ++(1,0.5); \node[font=\tiny\ttfamily,text=fg] at (34.5,-1.30) {0};
\draw[draw=acc!70!bg,fill=acc!45!bg] (35,-1.55) rectangle ++(1,0.5); \node[font=\tiny\ttfamily,text=fg] at (35.5,-1.30) {1};
\draw[draw=acc!70!bg,fill=acc!45!bg] (36,-1.55) rectangle ++(1,0.5); \node[font=\tiny\ttfamily,text=fg] at (36.5,-1.30) {1};
\draw[draw=acc!70!bg,fill=codebg] (37,-1.55) rectangle ++(1,0.5); \node[font=\tiny\ttfamily,text=fg] at (37.5,-1.30) {0};
\draw[draw=acc!70!bg,fill=codebg] (38,-1.55) rectangle ++(1,0.5); \node[font=\tiny\ttfamily,text=fg] at (38.5,-1.30) {0};
\draw[draw=acc!70!bg,fill=acc!45!bg] (39,-1.55) rectangle ++(1,0.5); \node[font=\tiny\ttfamily,text=fg] at (39.5,-1.30) {1};
\draw[draw=acc!70!bg,fill=acc!45!bg] (40,-1.55) rectangle ++(1,0.5); \node[font=\tiny\ttfamily,text=fg] at (40.5,-1.30) {1};
\draw[draw=acc!70!bg,fill=codebg] (41,-1.55) rectangle ++(1,0.5); \node[font=\tiny\ttfamily,text=fg] at (41.5,-1.30) {0};
\draw[draw=acc!70!bg,fill=codebg] (42,-1.55) rectangle ++(1,0.5); \node[font=\tiny\ttfamily,text=fg] at (42.5,-1.30) {0};
\draw[draw=acc!70!bg,fill=acc!45!bg] (43,-1.55) rectangle ++(1,0.5); \node[font=\tiny\ttfamily,text=fg] at (43.5,-1.30) {1};
\draw[draw=acc!70!bg,fill=acc!45!bg] (44,-1.55) rectangle ++(1,0.5); \node[font=\tiny\ttfamily,text=fg] at (44.5,-1.30) {1};
\draw[draw=acc!70!bg,fill=codebg] (45,-1.55) rectangle ++(1,0.5); \node[font=\tiny\ttfamily,text=fg] at (45.5,-1.30) {0};
\draw[draw=acc!70!bg,fill=codebg] (46,-1.55) rectangle ++(1,0.5); \node[font=\tiny\ttfamily,text=fg] at (46.5,-1.30) {0};
\draw[draw=acc!70!bg,fill=acc!45!bg] (47,-1.55) rectangle ++(1,0.5); \node[font=\tiny\ttfamily,text=fg] at (47.5,-1.30) {1};
\draw[draw=acc!70!bg,fill=acc!45!bg] (48,-1.55) rectangle ++(1,0.5); \node[font=\tiny\ttfamily,text=fg] at (48.5,-1.30) {1};
\draw[draw=acc!70!bg,fill=codebg] (49,-1.55) rectangle ++(1,0.5); \node[font=\tiny\ttfamily,text=fg] at (49.5,-1.30) {0};
\draw[draw=acc!70!bg,fill=codebg] (50,-1.55) rectangle ++(1,0.5); \node[font=\tiny\ttfamily,text=fg] at (50.5,-1.30) {0};
\draw[draw=acc!70!bg,fill=acc!45!bg] (51,-1.55) rectangle ++(1,0.5); \node[font=\tiny\ttfamily,text=fg] at (51.5,-1.30) {1};
\draw[draw=acc!70!bg,fill=acc!45!bg] (52,-1.55) rectangle ++(1,0.5); \node[font=\tiny\ttfamily,text=fg] at (52.5,-1.30) {1};
\draw[draw=acc!70!bg,fill=codebg] (53,-1.55) rectangle ++(1,0.5); \node[font=\tiny\ttfamily,text=fg] at (53.5,-1.30) {0};
\draw[draw=acc!70!bg,fill=codebg] (54,-1.55) rectangle ++(1,0.5); \node[font=\tiny\ttfamily,text=fg] at (54.5,-1.30) {0};
\draw[draw=acc!70!bg,fill=acc!45!bg] (55,-1.55) rectangle ++(1,0.5); \node[font=\tiny\ttfamily,text=fg] at (55.5,-1.30) {1};
\draw[draw=acc!70!bg,fill=acc!45!bg] (56,-1.55) rectangle ++(1,0.5); \node[font=\tiny\ttfamily,text=fg] at (56.5,-1.30) {1};
\draw[draw=acc!70!bg,fill=codebg] (57,-1.55) rectangle ++(1,0.5); \node[font=\tiny\ttfamily,text=fg] at (57.5,-1.30) {0};
\draw[draw=acc!70!bg,fill=codebg] (58,-1.55) rectangle ++(1,0.5); \node[font=\tiny\ttfamily,text=fg] at (58.5,-1.30) {0};
\draw[draw=acc!70!bg,fill=acc!45!bg] (59,-1.55) rectangle ++(1,0.5); \node[font=\tiny\ttfamily,text=fg] at (59.5,-1.30) {1};
\draw[draw=acc!70!bg,fill=acc!45!bg] (60,-1.55) rectangle ++(1,0.5); \node[font=\tiny\ttfamily,text=fg] at (60.5,-1.30) {1};
\draw[draw=acc!70!bg,fill=codebg] (61,-1.55) rectangle ++(1,0.5); \node[font=\tiny\ttfamily,text=fg] at (61.5,-1.30) {0};
\draw[draw=acc!70!bg,fill=acc!45!bg] (62,-1.55) rectangle ++(1,0.5); \node[font=\tiny\ttfamily,text=fg] at (62.5,-1.30) {1};
\draw[draw=acc!70!bg,fill=codebg] (63,-1.55) rectangle ++(1,0.5); \node[font=\tiny\ttfamily,text=fg] at (63.5,-1.30) {0};
\node[anchor=east,font=\small\ttfamily,text=fg] at (-0.6,-2.85) {1.0 - 0.9};
\node[anchor=west,font=\scriptsize\ttfamily,text=org] at (64.6,-2.85) {0x3FB9999999999998};
\draw[draw=red!70!bg,fill=codebg] (0,-3.1) rectangle ++(1,0.5); \node[font=\tiny\ttfamily,text=fg] at (0.5,-2.85) {0};
\draw[draw=yel!70!bg,fill=codebg] (1,-3.1) rectangle ++(1,0.5); \node[font=\tiny\ttfamily,text=fg] at (1.5,-2.85) {0};
\draw[draw=yel!70!bg,fill=yel!45!bg] (2,-3.1) rectangle ++(1,0.5); \node[font=\tiny\ttfamily,text=fg] at (2.5,-2.85) {1};
\draw[draw=yel!70!bg,fill=yel!45!bg] (3,-3.1) rectangle ++(1,0.5); \node[font=\tiny\ttfamily,text=fg] at (3.5,-2.85) {1};
\draw[draw=yel!70!bg,fill=yel!45!bg] (4,-3.1) rectangle ++(1,0.5); \node[font=\tiny\ttfamily,text=fg] at (4.5,-2.85) {1};
\draw[draw=yel!70!bg,fill=yel!45!bg] (5,-3.1) rectangle ++(1,0.5); \node[font=\tiny\ttfamily,text=fg] at (5.5,-2.85) {1};
\draw[draw=yel!70!bg,fill=yel!45!bg] (6,-3.1) rectangle ++(1,0.5); \node[font=\tiny\ttfamily,text=fg] at (6.5,-2.85) {1};
\draw[draw=yel!70!bg,fill=yel!45!bg] (7,-3.1) rectangle ++(1,0.5); \node[font=\tiny\ttfamily,text=fg] at (7.5,-2.85) {1};
\draw[draw=yel!70!bg,fill=yel!45!bg] (8,-3.1) rectangle ++(1,0.5); \node[font=\tiny\ttfamily,text=fg] at (8.5,-2.85) {1};
\draw[draw=yel!70!bg,fill=codebg] (9,-3.1) rectangle ++(1,0.5); \node[font=\tiny\ttfamily,text=fg] at (9.5,-2.85) {0};
\draw[draw=yel!70!bg,fill=yel!45!bg] (10,-3.1) rectangle ++(1,0.5); \node[font=\tiny\ttfamily,text=fg] at (10.5,-2.85) {1};
\draw[draw=yel!70!bg,fill=yel!45!bg] (11,-3.1) rectangle ++(1,0.5); \node[font=\tiny\ttfamily,text=fg] at (11.5,-2.85) {1};
\draw[draw=acc!70!bg,fill=acc!45!bg] (12,-3.1) rectangle ++(1,0.5); \node[font=\tiny\ttfamily,text=fg] at (12.5,-2.85) {1};
\draw[draw=acc!70!bg,fill=codebg] (13,-3.1) rectangle ++(1,0.5); \node[font=\tiny\ttfamily,text=fg] at (13.5,-2.85) {0};
\draw[draw=acc!70!bg,fill=codebg] (14,-3.1) rectangle ++(1,0.5); \node[font=\tiny\ttfamily,text=fg] at (14.5,-2.85) {0};
\draw[draw=acc!70!bg,fill=acc!45!bg] (15,-3.1) rectangle ++(1,0.5); \node[font=\tiny\ttfamily,text=fg] at (15.5,-2.85) {1};
\draw[draw=acc!70!bg,fill=acc!45!bg] (16,-3.1) rectangle ++(1,0.5); \node[font=\tiny\ttfamily,text=fg] at (16.5,-2.85) {1};
\draw[draw=acc!70!bg,fill=codebg] (17,-3.1) rectangle ++(1,0.5); \node[font=\tiny\ttfamily,text=fg] at (17.5,-2.85) {0};
\draw[draw=acc!70!bg,fill=codebg] (18,-3.1) rectangle ++(1,0.5); \node[font=\tiny\ttfamily,text=fg] at (18.5,-2.85) {0};
\draw[draw=acc!70!bg,fill=acc!45!bg] (19,-3.1) rectangle ++(1,0.5); \node[font=\tiny\ttfamily,text=fg] at (19.5,-2.85) {1};
\draw[draw=acc!70!bg,fill=acc!45!bg] (20,-3.1) rectangle ++(1,0.5); \node[font=\tiny\ttfamily,text=fg] at (20.5,-2.85) {1};
\draw[draw=acc!70!bg,fill=codebg] (21,-3.1) rectangle ++(1,0.5); \node[font=\tiny\ttfamily,text=fg] at (21.5,-2.85) {0};
\draw[draw=acc!70!bg,fill=codebg] (22,-3.1) rectangle ++(1,0.5); \node[font=\tiny\ttfamily,text=fg] at (22.5,-2.85) {0};
\draw[draw=acc!70!bg,fill=acc!45!bg] (23,-3.1) rectangle ++(1,0.5); \node[font=\tiny\ttfamily,text=fg] at (23.5,-2.85) {1};
\draw[draw=acc!70!bg,fill=acc!45!bg] (24,-3.1) rectangle ++(1,0.5); \node[font=\tiny\ttfamily,text=fg] at (24.5,-2.85) {1};
\draw[draw=acc!70!bg,fill=codebg] (25,-3.1) rectangle ++(1,0.5); \node[font=\tiny\ttfamily,text=fg] at (25.5,-2.85) {0};
\draw[draw=acc!70!bg,fill=codebg] (26,-3.1) rectangle ++(1,0.5); \node[font=\tiny\ttfamily,text=fg] at (26.5,-2.85) {0};
\draw[draw=acc!70!bg,fill=acc!45!bg] (27,-3.1) rectangle ++(1,0.5); \node[font=\tiny\ttfamily,text=fg] at (27.5,-2.85) {1};
\draw[draw=acc!70!bg,fill=acc!45!bg] (28,-3.1) rectangle ++(1,0.5); \node[font=\tiny\ttfamily,text=fg] at (28.5,-2.85) {1};
\draw[draw=acc!70!bg,fill=codebg] (29,-3.1) rectangle ++(1,0.5); \node[font=\tiny\ttfamily,text=fg] at (29.5,-2.85) {0};
\draw[draw=acc!70!bg,fill=codebg] (30,-3.1) rectangle ++(1,0.5); \node[font=\tiny\ttfamily,text=fg] at (30.5,-2.85) {0};
\draw[draw=acc!70!bg,fill=acc!45!bg] (31,-3.1) rectangle ++(1,0.5); \node[font=\tiny\ttfamily,text=fg] at (31.5,-2.85) {1};
\draw[draw=acc!70!bg,fill=acc!45!bg] (32,-3.1) rectangle ++(1,0.5); \node[font=\tiny\ttfamily,text=fg] at (32.5,-2.85) {1};
\draw[draw=acc!70!bg,fill=codebg] (33,-3.1) rectangle ++(1,0.5); \node[font=\tiny\ttfamily,text=fg] at (33.5,-2.85) {0};
\draw[draw=acc!70!bg,fill=codebg] (34,-3.1) rectangle ++(1,0.5); \node[font=\tiny\ttfamily,text=fg] at (34.5,-2.85) {0};
\draw[draw=acc!70!bg,fill=acc!45!bg] (35,-3.1) rectangle ++(1,0.5); \node[font=\tiny\ttfamily,text=fg] at (35.5,-2.85) {1};
\draw[draw=acc!70!bg,fill=acc!45!bg] (36,-3.1) rectangle ++(1,0.5); \node[font=\tiny\ttfamily,text=fg] at (36.5,-2.85) {1};
\draw[draw=acc!70!bg,fill=codebg] (37,-3.1) rectangle ++(1,0.5); \node[font=\tiny\ttfamily,text=fg] at (37.5,-2.85) {0};
\draw[draw=acc!70!bg,fill=codebg] (38,-3.1) rectangle ++(1,0.5); \node[font=\tiny\ttfamily,text=fg] at (38.5,-2.85) {0};
\draw[draw=acc!70!bg,fill=acc!45!bg] (39,-3.1) rectangle ++(1,0.5); \node[font=\tiny\ttfamily,text=fg] at (39.5,-2.85) {1};
\draw[draw=acc!70!bg,fill=acc!45!bg] (40,-3.1) rectangle ++(1,0.5); \node[font=\tiny\ttfamily,text=fg] at (40.5,-2.85) {1};
\draw[draw=acc!70!bg,fill=codebg] (41,-3.1) rectangle ++(1,0.5); \node[font=\tiny\ttfamily,text=fg] at (41.5,-2.85) {0};
\draw[draw=acc!70!bg,fill=codebg] (42,-3.1) rectangle ++(1,0.5); \node[font=\tiny\ttfamily,text=fg] at (42.5,-2.85) {0};
\draw[draw=acc!70!bg,fill=acc!45!bg] (43,-3.1) rectangle ++(1,0.5); \node[font=\tiny\ttfamily,text=fg] at (43.5,-2.85) {1};
\draw[draw=acc!70!bg,fill=acc!45!bg] (44,-3.1) rectangle ++(1,0.5); \node[font=\tiny\ttfamily,text=fg] at (44.5,-2.85) {1};
\draw[draw=acc!70!bg,fill=codebg] (45,-3.1) rectangle ++(1,0.5); \node[font=\tiny\ttfamily,text=fg] at (45.5,-2.85) {0};
\draw[draw=acc!70!bg,fill=codebg] (46,-3.1) rectangle ++(1,0.5); \node[font=\tiny\ttfamily,text=fg] at (46.5,-2.85) {0};
\draw[draw=acc!70!bg,fill=acc!45!bg] (47,-3.1) rectangle ++(1,0.5); \node[font=\tiny\ttfamily,text=fg] at (47.5,-2.85) {1};
\draw[draw=acc!70!bg,fill=acc!45!bg] (48,-3.1) rectangle ++(1,0.5); \node[font=\tiny\ttfamily,text=fg] at (48.5,-2.85) {1};
\draw[draw=acc!70!bg,fill=codebg] (49,-3.1) rectangle ++(1,0.5); \node[font=\tiny\ttfamily,text=fg] at (49.5,-2.85) {0};
\draw[draw=acc!70!bg,fill=codebg] (50,-3.1) rectangle ++(1,0.5); \node[font=\tiny\ttfamily,text=fg] at (50.5,-2.85) {0};
\draw[draw=acc!70!bg,fill=acc!45!bg] (51,-3.1) rectangle ++(1,0.5); \node[font=\tiny\ttfamily,text=fg] at (51.5,-2.85) {1};
\draw[draw=acc!70!bg,fill=acc!45!bg] (52,-3.1) rectangle ++(1,0.5); \node[font=\tiny\ttfamily,text=fg] at (52.5,-2.85) {1};
\draw[draw=acc!70!bg,fill=codebg] (53,-3.1) rectangle ++(1,0.5); \node[font=\tiny\ttfamily,text=fg] at (53.5,-2.85) {0};
\draw[draw=acc!70!bg,fill=codebg] (54,-3.1) rectangle ++(1,0.5); \node[font=\tiny\ttfamily,text=fg] at (54.5,-2.85) {0};
\draw[draw=acc!70!bg,fill=acc!45!bg] (55,-3.1) rectangle ++(1,0.5); \node[font=\tiny\ttfamily,text=fg] at (55.5,-2.85) {1};
\draw[draw=acc!70!bg,fill=acc!45!bg] (56,-3.1) rectangle ++(1,0.5); \node[font=\tiny\ttfamily,text=fg] at (56.5,-2.85) {1};
\draw[draw=acc!70!bg,fill=codebg] (57,-3.1) rectangle ++(1,0.5); \node[font=\tiny\ttfamily,text=fg] at (57.5,-2.85) {0};
\draw[draw=acc!70!bg,fill=codebg] (58,-3.1) rectangle ++(1,0.5); \node[font=\tiny\ttfamily,text=fg] at (58.5,-2.85) {0};
\draw[draw=acc!70!bg,fill=acc!45!bg] (59,-3.1) rectangle ++(1,0.5); \node[font=\tiny\ttfamily,text=fg] at (59.5,-2.85) {1};
\draw[draw=acc!70!bg,fill=acc!45!bg] (60,-3.1) rectangle ++(1,0.5); \node[font=\tiny\ttfamily,text=fg] at (60.5,-2.85) {1};
\draw[draw=acc!70!bg,fill=codebg] (61,-3.1) rectangle ++(1,0.5); \node[font=\tiny\ttfamily,text=fg] at (61.5,-2.85) {0};
\draw[draw=acc!70!bg,fill=codebg] (62,-3.1) rectangle ++(1,0.5); \node[font=\tiny\ttfamily,text=fg] at (62.5,-2.85) {0};
\draw[draw=acc!70!bg,fill=codebg] (63,-3.1) rectangle ++(1,0.5); \node[font=\tiny\ttfamily,text=fg] at (63.5,-2.85) {0};
\draw[decorate,decoration={brace,amplitude=4pt},draw=red] (0,0.65) -- (1,0.65) node[midway,above left=5pt and -2pt,font=\scriptsize,text=red]{sign};
\draw[decorate,decoration={brace,amplitude=4pt},draw=yel] (1,0.65) -- (12,0.65) node[midway,above=4pt,xshift=14pt,font=\scriptsize,text=yel]{exponent (11 bits, bias 1023)};
\draw[decorate,decoration={brace,amplitude=4pt},draw=acc] (12,0.65) -- (64,0.65) node[midway,above=4pt,font=\scriptsize,text=acc]{fraction (52 bits; leading 1 implicit)};
\end{tikzpicture}
}
\caption{The 64 bits of \texttt{0.9}, \texttt{0.1} and \texttt{1.0 - 0.9} (TikZ, generated from \cc{struct.pack} by
\file{figures/gen_bits.py}). Filled cells are 1s. The repeating \texttt{1100} / \texttt{1001} pattern is $1/10$'s
infinite binary expansion, cut after 52 fraction bits and rounded; the last two hex digits of the third row,
\texttt{98}, differ from the second row's \texttt{9A}.}
\label{fig:bits}
\end{figure}

\begin{verified}
From \file{verify/fp.py} (Python \cc{struct}, \cc{decimal.Decimal}, \cc{fractions.Fraction}):
\begin{center}
\footnotesize
\begin{tabular}{@{}lll@{}}
\toprule
\textcolor{muted}{literal} & \textcolor{muted}{bits (hex)} & \textcolor{muted}{exact stored value}\\
\midrule
\texttt{0.9} & \hex{0x3FECCCCCCCCCCCCD} & $0.90000000000000002220446049250313080847263336181640625$\\
\texttt{0.1} & \hex{0x3FB999999999999A} & $0.1000000000000000055511151231257827021181583404541015625$\\
\texttt{1.0-0.9} & \hex{0x3FB9999999999998} & $0.09999999999999997779553950749686919152736663818359375$\\
\texttt{0.95} & \hex{0x3FEE666666666666} & $0.9499999999999999555910790149937383830547332763671875$\\
\texttt{0.8} & \hex{0x3FE999999999999A} & $0.8000000000000000444089209850062616169452667236328125$\\
\bottomrule
\end{tabular}
\end{center}
As fractions, $\fl(0.9)=8106479329266893/2^{53}$ and $\fl(1.0-0.9)=900719925474099/2^{53}$.
\end{verified}

So $\fl(0.9)$ is slightly \emph{above} $0.9$ (by $2.22\times10^{-17}$), and therefore $1-\fl(0.9)$ is slightly
\emph{below} $0.1$. Printed with the 17 significant digits needed to identify a binary64 uniquely, it is
\texttt{0.09999999999999998}. Conversely $\fl(0.95)$ is below $0.95$, so $1-\fl(0.95)$ is \emph{above} $0.05$: that
rounding direction happened to be harmless.

\section{Where the error comes from: Sterbenz's lemma}
It is tempting to blame the subtraction. It is innocent.

\begin{lemma}[Sterbenz]\label[lemma]{lem:sterbenz}
If $x$ and $y$ are binary64 numbers with $y/2\le x\le 2y$, then $x-y$ is exactly representable, so
$\fl(x-y)=x-y$.
\end{lemma}
\begin{proof}[Proof sketch]
Assume $0<y\le x\le 2y$ (the other case is symmetric). Both are integer multiples of $\ulp(y)$, the spacing at $y$, because
$x$ has an exponent at least that of $y$. So $x-y$ is a multiple of $\ulp(y)$, and $0\le x-y\le y$, so it needs no more
significant bits than $y$ has. Hence it is representable. (Goldberg gives the full argument, including subnormals~\cite{goldberg1991what}.)
\end{proof}

With $x=1$ and $y=\fl(0.9)$, $y/2\le1\le2y$ holds, so \cc{1.0 - 0.9} is computed exactly: the whole discrepancy is the
representation error of the \emph{literal} $0.9$. The multiplication by $10$ is exact too in this case (verified in Python:
\cc{Fraction((1.0-0.9)*10) == Fraction(1.0-0.9)*10} is \cc{True}), giving
\[ 0.99999999999999977795539507496869191527366638183593750, \]
which \cc{static_cast<std::size_t>} truncates to $0$.

\begin{layman}
Imagine a ruler marked only in sixteenths of an inch, and you are told to measure nine tenths of an inch. You pick the
nearest mark, slightly too long. ``One inch minus that'' is then measured perfectly, and it is slightly short of a tenth,
because the error was in the mark you chose, not in the subtraction. Ten of those slightly-short tenths do not reach the
one-inch line, and ``how many whole inches?'' answers zero.
\end{layman}

\section{When is \texorpdfstring{$\varepsilon=10^{-9}$}{epsilon = 1e-9} enough?}
We model the computation exactly as the code performs it, using the standard model
$\fl(a\circ b)=(a\circ b)(1+\delta)$, $|\delta|\le u=2^{-53}$~\cite{goldberg1991what,higham2002accuracy}.

\begin{theorem}[Correctness of the epsilon floor]\label{thm:epsfloor}
Let the confidence level be a decimal with at most $d$ digits after the point, $c=m/10^d\in(0,1)$, let
$1\le n<2^{53}$, and let the code compute $\hat k=\bigl\lfloor\fl\bigl(\fl(\fl(1-\hat c)\cdot n)+\varepsilon\bigr)\bigr\rfloor$
with $\hat c=\fl(c)$. Put $\eta(n)=n\cdot2^{-51}$. If
\[ \eta(n)<\varepsilon \qquad\text{and}\qquad \varepsilon+\eta(n)<10^{-d}, \]
then $\hat k=\lfloor(1-c)\,n\rfloor$ exactly.
\end{theorem}

\begin{proof}
Let $y=(1-c)n$ be the exact value and $K=\lfloor y\rfloor$.

\emph{Step 1: the tail fraction.} Since $c<1$, $|\hat c-c|\le\frac12\ulp(c)\le2^{-54}$. If $\hat c\ge\frac12$, Sterbenz
(\cref{lem:sterbenz}) makes $1-\hat c$ exact; otherwise $1-\hat c\in(\frac12,1]$ and rounding adds at most $2^{-54}$. In
both cases $t:=\fl(1-\hat c)$ satisfies $|t-(1-c)|\le2^{-53}$.

\emph{Step 2: the product.} $\hat y:=\fl(tn)$ satisfies $|\hat y-tn|\le u\,tn\le n2^{-53}$, so
$|\hat y-y|\le|\hat y-tn|+n|t-(1-c)|\le n2^{-52}$.

\emph{Step 3: adding $\varepsilon$.} $z:=\fl(\hat y+\varepsilon)$ satisfies $|z-(\hat y+\varepsilon)|\le u(\hat y+\varepsilon)\le n2^{-52}$
for $n\ge1$. Altogether $|z-(y+\varepsilon)|\le n2^{-51}=\eta(n)$.

\emph{Step 4: the floor.} Because $y=(10^d-m)n/10^d$, its fractional part $\phi=y-K$ is a multiple of $10^{-d}$, so
$\phi\le1-10^{-d}$. Then
\[ z\ \ge\ y+\varepsilon-\eta\ \ge\ K+(\varepsilon-\eta)\ >\ K,\qquad
   z\ \le\ K+\phi+\varepsilon+\eta\ \le\ K+1-10^{-d}+\varepsilon+\eta\ <\ K+1, \]
by the two hypotheses. Hence $\lfloor z\rfloor=K$; for non-negative $z$ the cast to \cc{std::size_t} truncates, which is the floor.
\end{proof}

\begin{corollary}[What the repository's constant buys]\label[corollary]{cor:eps}
With $\varepsilon=10^{-9}$ the hypotheses hold for every confidence level with at most $8$ decimal digits and every
$n\le2{,}251{,}799$ ($=\lfloor10^{-9}\cdot2^{51}\rfloor$). With $\varepsilon=0$ the first hypothesis fails, and the floor can be wrong
whenever $y$ is an integer and the accumulated rounding error is negative, as at $(n,c)=(10,0.90)$.
\end{corollary}

The bound on $n$ is about 8{,}900 years of daily data, so it does not bind. The bound is also loose; experiment shows more
headroom.

\begin{verified}
\file{verify/fp.py} compares the guarded and unguarded expressions with $\lfloor(1-c)n\rfloor$ computed in exact rational
arithmetic for all $9{,}999$ four-decimal levels $c\in\{0.0001,\dots,0.9999\}$ and all
$n\in\{1,\dots,2000\}\cup\{10^4,10^5,10^6\}$: $20{,}027{,}997$ cases. The guarded expression has \textbf{0} mismatches;
the unguarded one has \textbf{9{,}344}.
\end{verified}

\subsection{Where the guard fails}
The second hypothesis is not decoration. Take a level with eleven decimals, $c=0.90000000005$, and $n=10$. Exactly,
$(1-c)\cdot10=0.9999999995$, so the true tail size is $0$ and the function should throw. The computed product is
$0.9999999994999997$; adding $10^{-9}$ pushes it past $1$.

\begin{progout}{eps\_fail: the repository's historical\_var with c = 0.90000000005, n = 10}{lst:epsfail}
c=0.90000000005000003 -> VaR=0.050000000000000003 (true tail size floor((1-c)*10)=0, should throw)
msg for c=0.9500001: risk::historical_var: not enough returns (10) to resolve the 0.950000 confidence quantile
\end{progout}

The second line shows a smaller wart: the error message formats the level with \cc{std::to_string}, which prints six
decimals, so a caller who passed $0.9500001$ is told about ``0.950000''.

\section{Alternatives to an epsilon}
\begin{center}
\small
\begin{tabularx}{\linewidth}{@{}>{\leavevmode\bfseries\color{yel}}p{2.9cm} X X@{}}
\toprule
\textcolor{muted}{Approach} & \textcolor{muted}{How} & \textcolor{muted}{Trade-off}\\
\midrule
$\varepsilon$-floor (current) & $\lfloor\fl((1-c)n)+10^{-9}\rfloor$ & One line; correct under \cref{thm:epsfloor}; wrong for levels with more than eight decimals.\\
Integer basis points & Take \cc{conf_bp} $\in[1,9999]$; $k=\lfloor(10^4-\text{bp})\,n/10^4\rfloor$ in \cc{uint64_t} & Exact for $n<1.8\times10^{15}$; changes the API; cannot express $c=0.99995$.\\
Round with tolerance & $k=\mathrm{round}(pn)$ if $|pn-\mathrm{round}(pn)|<\tau$, else $\lfloor pn\rfloor$ & Same idea as $\varepsilon$, stated symmetrically; still needs a $\tau$ and a theorem.\\
Exact rationals & Parse the level as a decimal string into a fraction; compute $\lfloor\cdot\rfloor$ exactly & Always right; the confidence must arrive as text (it does, over the JSON wire).\\
Decimal floating point & IEEE-754 \texttt{decimal64}~\cite{ieee754-2019} & $0.9$ is exact; little native hardware or standard-library support in C++.\\
\bottomrule
\end{tabularx}
\end{center}

\begin{code}{tail\_index\_bp: confidence in basis points (not repository code)}{lst:bp}
std::size_t tail_index_bp(std::uint32_t conf_bp, std::size_t n) {
    if (conf_bp == 0 || conf_bp >= 10'000) throw std::domain_error("conf_bp must be in [1, 9999]");
    // exact: floor((1 - conf_bp/10^4) * n) == floor((10^4 - conf_bp) * n / 10^4)
    return static_cast<std::size_t>((std::uint64_t{10'000 - conf_bp} * n) / 10'000);
}
\end{code}

\begin{verified}
\file{verify/bp.cpp} checks \cref{lst:bp} against a brute-force integer floor for all $9{,}999$ levels and $n\le61$
($609{,}939$ cases, 0 mismatches), and reproduces $k=1,0,2$ at $n=10$ for 9000, 9500 and 8000 basis points and
$k=25,12,5$ at $n=503$ for 9500, 9750 and 9900. The product overflows \cc{uint64_t} only for
$n>\lfloor2^{64}/10^4\rfloor=1{,}844{,}674{,}407{,}370{,}955$.
\end{verified}

\begin{remarkbox}
The wire format already carries the confidence as text: \cref{fig:frame} shows \texttt{0.95} as the bytes
\hex{30 2E 39 35}. A server could parse that decimal string straight into basis points or a rational and never create the
inexact \cc{double} at all. Choosing \cc{double} for a quantity that users think of as an exact decimal is the root cause;
the epsilon treats the symptom, provably well within the stated bounds.
\end{remarkbox}

\begin{history}
Three well-documented failures share this chapter's anatomy: a quantity that people think of as exact, stored in a
format that cannot hold it, with the error accumulating or overflowing in a place nobody was looking. The Patriot
battery at Dhahran (1991) kept time in tenths of a second with a truncated binary $1/10$; the Vancouver Stock Exchange index
(1982--83) truncated instead of rounding at every recomputation; Ariane~5 Flight~501 (1996) converted a 64-bit float to a
16-bit integer that could not hold it. \Cref{sec:hist-patriot,sec:hist-vse,sec:hist-ariane} work through each with numbers we recomputed.
\end{history}

\begin{exercises}
\ex{ex:095}{\lvA} Without running code, predict from \cref{fig:bits}'s pattern whether $\fl(0.95)$ lies above or below
$0.95$ (hint: look at the last fraction bits of $19/20$'s expansion and the rounding). Confirm with the table above, and
explain why \cc{1.0 - 0.95} is greater than $0.05$.

\ex{ex:sterbenz}{\lvB} Complete the proof of \cref{lem:sterbenz} for normal numbers, including the case where $x$ and $y$
have different exponents. Then show the lemma fails for $x=1$, $y=\fl(0.1)$ by finding the rounding error of \cc{1.0 - 0.1}.

\ex{ex:characterise}{\lvB} For $c\in\{0.01,0.02,\dots,0.99\}$ and $n\le1000$, find every pair where the \emph{unguarded}
floor is wrong. Characterise them: is it always the case that $(1-c)n$ is an integer? Can the unguarded floor ever be too
\emph{large}? Prove your answer.

\ex{ex:bpapi}{\lvB} Change \cc{historical_var} to compute $k$ by \cref{lst:bp} internally while keeping the
\cc{double confidence_level} signature for Python callers. How do you convert the \cc{double} to basis points safely, and
which inputs do you now reject that were accepted before?

\ex{ex:message}{\lvA} Fix the error message so that it prints the confidence level exactly as the shortest decimal that
round-trips (C++23 \cc{std::format("{}", x)} does this). Produce one input for which the old and new messages differ.

\ex{ex:kahan}{\lvB} \cc{historical_cvar} sums $k$ sorted returns with \cc{std::accumulate}. Bound the absolute error of the
sum with the standard model, evaluate the bound for the AAPL mock series at $k=25$, and decide whether compensated (Kahan)
summation could change any of the six significant digits reported in \cref{tab:mockvar}.
\end{exercises}

